% 2.3 =============== 值函数性质：数学准备 ===============

\section{数学准备} \label{sec-math-pre}

上面的例子其实引出了两个问题：
\begin{itemize}
    \item 通过定义贝尔曼算子，我们将递归问题变为不动点问题$v=Tv$的求解，但是这个解在理论上是否存在和唯一？
    \item 给定初始猜测的$v_0\left(k\right)$，由算子$T$生成的值函数序列$\{v_n\}$是否收敛到$v$？即这一算法是否可行？
\end{itemize}

证明这两个问题需要做一些泛函分析的数学准备,例如，如何定义函数间的距离、函数序列的收敛性等。


% 2.3.1 ===== 度量空间 =====
\subsection{度量空间中的收敛性}

\begin{definition}
    \textbf{度量空间}（Metric Space）由一个非空集合$S$和某一度量或称为距离（distance function
    ）$d: S \times S \rightarrow \mathbb{R}$所定义，其中度量$d$满足：$\forall x,y,z \in S$
    \begin{itemize}
        \item $d\left(x,y\right) \geq 0$，当且仅当$x=y$时取等
        \item $d\left(x,y\right) = d\left(y,x\right)$
        \item $d(x,y) \leq d\left(x,z\right)+d\left(z,y\right) $
    \end{itemize}
\end{definition}

在这里，集合$S$不止于数集，也可能是函数集（例如我们感兴趣的值函数集）等。
因此，度量空间的定义比普通欧氏距离（Euclidean Distance）更复杂。
但显然，实数集$S = \mathbb{R}$和绝对距离$d\left(x,y\right)=\left|x-y\right|$构成一个度量空间。
\begin{example}
    常见的度量空间例子：
    \begin{itemize}
        \item 任意有限维空间$\mathbb{R}^n$和其上定义的距离：
            \begin{gather*}
                d\left(x,y\right) = \left(\sum_{i=1}^{n} \left|x_i-y_i\right|^p\right)^{1/p},\quad 1 \leq p < \infty \\
                d\left(x,y\right) = \max _i \left(\left|x_i-y_i\right|\right), \quad p=\infty
            \end{gather*}

        \item 定义在$\left[a,b\right]$上的连续函数集（常记为$C\left[a,b\right]$）和其上定义的距离：     
            \begin{gather*}
                d\left(x,y\right)=\left(\int_{a}^{b} \left|x\left(t\right)-y\left(t\right)\right|^{p}dt\right)^{1/p}, 
                    \quad 1 \leq p < \infty \\
                d\left(x,y\right) = \max_{a\leq t \leq b} \left(\left|x\left(t\right)-y\left(t\right)\right|\right), 
                    \quad p = \infty
            \end{gather*}
    \end{itemize}
\end{example}

\begin{definition}
    \textbf{赋范向量空间}（Normed Vector Space）由一个向量空间$S$和范数$\| \cdot \|: S \rightarrow \mathbb{R}$所定义，
    其中范数满足：$\forall x,y \in S$和$a \in \mathbb{R}$
    \begin{itemize}
        \item \( \|x\| = 0 \)当且仅当 \(x=\mathbf{0}\)
        \item $\|x\|=|a|\cdot \|x\|$
        \item $\|x+y\| \leq \|x\|+\|y\|$
    \end{itemize}
\end{definition}

不难看出，若令距离$d\left(x,y\right)=\|x-y\|$，那么任意赋范向量空间都是度量空间。
\begin{example}
    常见的赋范向量空间例子：
    \begin{itemize}
        \item 任意有限维空间$\mathbb{R}^n$和其上定义的范数：
            \begin{gather*}
                \|x\|=\left(\sum_{i=1}^{n}|x_i|^p\right)^{1/p}, \quad 1\leq p < \infty \\
                \|x\|=\max_{i}\left(|x_i|\right), \quad p=\infty
            \end{gather*}
        \item 定义在$\left[a,b\right]$上的连续函数集$C\left[a,b\right]$和其上定义的范数：     
            \begin{gather*}
                \|x\|=\left(\int_{a}^{b}|x\left(t\right)|^{p}dt\right)^{1/p}, \quad 1 \leq p < \infty \\
                \|x\|=\max _{a\leq t \leq b}\left(|x\left(t\right)|\right), \quad p = \infty
            \end{gather*}
    \end{itemize}
\end{example}

\begin{definition}[\textbf{收敛性}]
    我们称$S$中的序列$\{x_n\}_{n=0}^{\infty}$收敛到$x\in S$：若对$\forall \epsilon > 0$，存在$N_{\epsilon}$，使得：
    \begin{equation*}
        d\left(x_n,x\right)<\epsilon, \quad n\geq N_{\epsilon}
    \end{equation*}
\end{definition}

注意序列$\{x_n\}$不只限于实数（事实上在动态规划中，我们感兴趣的是函数序列）。但从另一角度看，
序列$\{x_n\}$收敛到$x \in S$等价于由距离构成的的实数序列$\{d\left(x_n,x\right)\}$收敛到0。
\begin{definition}[\textbf{柯西序列}]
    $S$中的序列$\{x_n\}_{n=0}^{\infty}$被称为柯西序列（Cauchy Sequence）若其满足柯西准测：
    对$\forall \epsilon > 0$，存在$N_{\epsilon}$，使得：
    \begin{equation*}
        d\left(x_n,x_m\right) < \epsilon, \quad n,m \geq N_{\epsilon}
    \end{equation*}
\end{definition}
\begin{theorem}
    所有收敛序列都是柯西序列。
\end{theorem}
\begin{theorem}
    所有柯西序列都有界。
\end{theorem}

定义柯西序列的好处在于，只利用序列本身信息就可以进行判断收敛性，而不用知道实际极限。
但是需要注意，\textbf{并非所有柯西序列都收敛}。实际上，这与序列所处空间的\textbf{完备性}（Completeness）密切相关。
\begin{example}
    考虑$\left[0,1\right]$上的连续函数集合$C\left[0,1\right]$。定义$L^1$度量：
    \begin{equation*}
        d\left(f,g\right) = \int_{0}^{1}\left|f\left(x\right)-g\left(x\right)\right|dx
    \end{equation*}
    构造函数序列$\{f_n\}$：
    \begin{equation*}
    f_n(x) = 
        \begin{cases}0 & x \in\left[0, \frac{1}{2}-\frac{1}{n}\right] \\ 
            n\left(x-\left(\frac{1}{2}-\frac{1}{n}\right)\right) & x 
                \in\left(\frac{1}{2}-\frac{1}{n}, \frac{1}{2}\right] \\ 
            1 & x \in\left(\frac{1}{2}, 1\right] 
        \end{cases}
    \end{equation*}
    由于对任意$m>n$，成立：
    \begin{equation*}
        d\left(f_m, f_n\right)=\int_0^1\left|f_m(x)-f_n(x)\right| d x \leq \frac{1}{n} \rightarrow 0 
            \quad(n \rightarrow \infty)
    \end{equation*}
    $\{f_n\}$在$L^1$度量下是柯西序列，但其收敛到极限为：
    \begin{equation*}
        f(x) = 
            \begin{cases}0 & x \in\left[0, \frac{1}{2}\right] \\ 
                1 & x \in\left(\frac{1}{2}, 1\right]
            \end{cases}
    \end{equation*}
    在$x=\frac{1}{2}$处不连续，因此极限函数$f \notin C[0,1]$。从而 $\left\{f_n\right\}$ 在 $C[0,1]$ 中无极限，不是收敛序列。
\end{example}

\begin{definition}[\textbf{完备性}]
    如果$S$中任意柯西序列都收敛到$S$中某一极限，则称度量空间$\left(S,d\right)$是完备的（Complete）。
    完备的赋范向量空间被称为巴拿赫（Banach）空间。
\end{definition}

\begin{theorem} \label{th-complete-continuous-func-space}
    令$X \subseteq \mathbb{R}^n$，定义了上确界范数：$\|f\|=\sup _{x \in X}(|f(x)|)$的所有有界连续函数
    $f: X \rightarrow \mathbb{R}$的集合$C(X)$是一个完备空间。
\end{theorem}

定义完备空间的方便之处在于，只要其中的序列符合柯西准测，那么它在这一空间中就存在极限。


% 2.2.2 ===== CMT =====
\subsection{压缩映射定理}

我们已经指出，贝尔曼方程\cref{eq-rck-bellman}的本质是求解一个泛函中的不动点问题：$v=Tv$。其中$T$是贝尔曼算子：
\begin{equation*}
    Tv\left(k\right) \equiv \max_{k^{\prime}} 
        \left[u\left(\gamma\left(k\right)-k^{\prime}\right)+\beta v\left(k^{\prime}\right)\right]
\end{equation*}
同时,从算法层面（寻找不动点）我们已经构造了递增的值函数序列$\{v_n\}$：
\begin{equation*}
    \begin{aligned}
        Tv_n\left(k\right) &\equiv \max _{k^{\prime}} 
            \left[u\left(\gamma\left(k\right)-k^{\prime}\right)+\beta v_n\left(k^{\prime}\right)\right] \\
        v_{n+1} &= Tv_n \geq v_n
    \end{aligned}
\end{equation*}
我们希望证明，随着迭代次数$n$足够大，值函数序列$\{v_n\}$收敛到唯一极限$v$。这需要用到\textbf{压缩映射定理}
（Contraction Mapping Theorem）。什么是压缩映射？和前面定义的度量空间和完备性有什么关系？

\begin{definition}[\textbf{压缩映射}]
    令$\left(S,d\right)$为度量空间，并定义$T:S\rightarrow S$。如果对某一$\beta \in \left(0,1\right)$，成立：
    \begin{equation*}
        d\left(Tx,Ty\right) \leq \beta d\left(x,y\right)
    \end{equation*}
    则称$T$为压缩映射，$\beta$为压缩系数。
\end{definition}
\begin{example}
    令$S=\left[a,b\right]$，$d\left(x,y\right)=\left|x-y\right|$。当对某一$\beta \in \left(0,1\right)$：
    \begin{equation*}
        \left|Tx-Ty\right|\leq \beta \left|x-y\right|
    \end{equation*}
    则$T:S\rightarrow S$为压缩映射。
    考虑$a=0,b=1$，$Tx=0.2+0.6x$：
    \begin{equation*}
         \left|Tx-Ty\right|=0.6\left(x-y\right)\leq \beta\left|x-y\right|,\quad \forall \beta \in [0.6,1)
    \end{equation*}
    因此$T$是压缩映射。
\end{example}

\begin{definition}
    若$x \in S$满足：$Tx=x$，则称$x$为不动点。
\end{definition}

\begin{theorem}[\textbf{压缩映射定理}]   \label{th-cmt}
    令$\left(S,d\right)$为完备度量空间，令$T:S\rightarrow S$为压缩映射，压缩系数为$\beta$，则
    \begin{itemize}
        \item $T$在$S$中有唯一不动点；
        \item 对任意$x_0 \in S$，成立：
            $d\left(T^nx_0, x\right) \leq \beta^n d\left(x_0,x\right), \quad n=0,1,2,\dots$
    \end{itemize}
\end{theorem}
\begin{proof}
    先证明不动点存在，接着利用压缩映射定义可得第二条。前者分三步证明：

    \textbf{Step1：证明迭代算子$T$后形成的序列是一个柯西序列}

    任取一个初始$x_0 \in S$，令$x_n=T^n x_0$，进而得到了序列$\{x_n\}_{n=0}^{\infty}$并且满足$x_{n+1}=Tx_n$。
    由于$T$是压缩映射，因此有：$d\left(x_2, x_1\right)=d\left(T x_1, T x_0\right) \leq \beta d\left(x_1, x_0\right)$,
    以此类推，对$n\geq 1$成立：
    \begin{equation*}
        d\left(x_{n+1}, x_n\right) \leq \beta^n d\left(x_1, x_0\right)
    \end{equation*}
    因此对任意$m>n$，应用距离的三角不等式性质：
    \begin{equation*}
        \begin{aligned}
            d\left(x_m, x_n\right) & \leq d\left(x_m, x_{m-1}\right) + 
                \cdots+d\left(x_{n+2}, x_{n+1}\right)+d\left(x_{n+1}, x_n\right) \\
            & \leq\left[\beta^{m-1}+\cdots+\beta^{n+1}+\beta^n\right] d\left(x_1, x_0\right) \\
            & =\beta^n\left[\beta^{m-n-1}+\cdots+\beta+1\right] d\left(x_1, x_0\right) \\
            & \leq \frac{\beta^n}{1-\beta} d\left(x_1, x_0\right)
        \end{aligned}
    \end{equation*}
    由于$\beta \in \left(0,1\right)$，显然$\{x_n\}$为柯西序列。又由于$S$是完备的，因此$x_n\rightarrow x \in S$

    \textbf{Step2：证明极限是$T$的不动点}

    再次运用不三角不等式和压缩映射的定义：
    \begin{equation*}
        \begin{aligned}
            d(T x, x) & \leq d\left(T x, T^n x_0\right)+d\left(T^n x_0, x\right) \\
            &= d\left(Tx,T\left(T^{n-1}x_0\right)\right) + \left(T^n x_0, x\right) \\
            & \leq \beta d\left(x, T^{n-1} x_0\right)+d\left(T^n x_0, x\right)
        \end{aligned} 
    \end{equation*}
    显然当$n\rightarrow \infty$，右边两项的距离都趋于0，因此：
    $$d\left(Tx,x\right)=0, \quad Tx = x$$
    从而序列$\{x_n\}$的极限$x$是$T$的不动点。

    \textbf{Step3：证明不动点的唯一性}

    反证，假设存在另一不动点$Tx^{\prime}=x^{\prime},x^{\prime} \neq x$，则：
    \begin{equation*}
        0<\delta \equiv d\left(x,x^{\prime}\right) = 
            d\left(Tx,Tx^{\prime}\right)\leq \beta d\left(x,x^{\prime}\right) =\beta \delta
    \end{equation*}
    但由于$\beta \in \left(0,1\right)$，显然矛盾，从而唯一性得证。

    最后，由压缩映射的定义立得：
    \begin{equation*}
        d\left(T^nx_0,x\right) = d\left(T\left(T^{n-1}x_0\right),Tx\right) \leq \beta d\left(T^{n-1}x_0,x\right)
    \end{equation*}
\end{proof}

压缩映射定理（Contraction Mapping Theorem）表明：
（1）\textbf{完备空间}上定义的压缩映射存在唯一不动点；（2）从\textbf{完备空间}$S$中任意初始点$x_0 \in S$出发，
在其上迭代压缩映射$T$，可以无限接近唯一的不动点。即我们找到了一个求解不动点的算法。

\textbf{再次回顾目标}：我们通过迭代贝尔曼算子：
\begin{equation*}
    Tv\left(k\right) \equiv \max_{k^{\prime}} 
        \left[u\left(\gamma\left(k\right)-k^{\prime}\right)+\beta v\left(k^{\prime}\right)\right]
\end{equation*}
得到了递增的值函数序列$v_{n+1}=Tv_n \geq v_n$，我们希望这一序列收敛到贝尔曼方程的不动点$v=Tv$。
我们已经指出，首先需要证明贝尔曼方程（不动点问题）确实存在唯一解；
其次去证明递增的值函数序列$\{v_n\}$收敛到不动点。显然，\textbf{如果能够证明贝尔曼算子是压缩映射}，那么直接应用压缩映射定理即可。
但是利用定义证明贝尔曼算子是压缩映射的难度较大，在这里给出一个充分条件：
\begin{theorem}[\textbf{Blackwell充分条件}]  \label{th-blackwell}
    令$X \subseteq \mathbb{R}^n$，令$B(X)$为有界函数集合$f: X \rightarrow \mathbb{R}$，
    其上定义了上确界范数（sup norm）：$\|f\|=\sup_{x \in X}(|f(x)|)$。
    令$T:B\left(X\right)\rightarrow B\left(X\right)$，且满足：
    \begin{itemize}
        \item 单调性：对$\forall f, g \in B(X)$，$f \leq g$， 则$T f \leq T g$
        \item 折扣性：存在$\beta \in \left(0,1\right)$，使得：$T(f+a) \leq T f+\beta a, \quad \forall f \in B(X)$。
            其中，函数$\left(f+a\right)\equiv f\left(x\right)+a$，$a \geq 0$。
    \end{itemize}
    那么$T$为压缩映射。
\end{theorem}
\begin{proof}
    令$v, w \in B(X)$，从而：
    \begin{equation*}
        \begin{aligned}
            v(x) & =v(x)+w(x)-w(x) \\
            & \leq w(x)+|v(x)-w(x)| \\
            & \leq w(x)+\|v-w\|, \quad \text { for all } x \in X
        \end{aligned}
    \end{equation*}
    由于映射$T$满足单调性和折扣性，因此：
    \begin{equation*}
        T v \leq T(w+\|v-w\|) \leq T w+\beta\|v-w\|
    \end{equation*}
    同理：
    \begin{equation*}
        T w \leq T(v+\|v-w\|) \leq T v+\beta\|v-w\|
    \end{equation*}
    根据这两个不等式可知：
    \begin{equation*}
        |T v(x)-T w(x)| \leq \beta\|v-w\|, \quad \text { for all } x \in X
    \end{equation*}
    也即：
    \begin{equation*}
        \|T v-T w\| \leq \beta\|v-w\|
    \end{equation*}
    从而$T$是压缩映射。
\end{proof}


% 2.3.3 ===== 应用CMT =====
\subsection{应用CMT}

重新考虑递归形式的最优增长模型。我们定义了贝尔曼方程右边为贝尔曼算子：
\begin{equation*}
    Tv\left(k\right) \equiv \max_{k^{\prime}} \left[u\left(\gamma\left(k\right)-k^{\prime}\right)
        +\beta v\left(k^{\prime}\right)\right]
\end{equation*}
下面验证这一算子的性质.

（1）\textbf{$T$满足单调性：}

若$v \leq w$，则:
\begin{equation*}
    u\left(\gamma\left(k\right)-k^{\prime}\right)+\beta v\left(k^{\prime}\right) \leq 
        u\left(\gamma\left(k\right)-k^{\prime}\right)+\beta w\left(k^{\prime}\right), \quad \forall x
\end{equation*}
从而有：
\begin{equation*}
    \max _{k^{\prime}}\left[u\left(\gamma\left(k\right)-k^{\prime}\right)+\beta v\left(k^{\prime}\right)\right] \leq 
    \max _{k^{\prime}}\left[u\left(\gamma\left(k\right)-k^{\prime}\right)+\beta w\left(k^{\prime}\right)\right]
\end{equation*}
因此：
$$Tv\left(k\right)\leq Tw\left(k\right)$$

（2）\textbf{T满足discounting：}

对$\forall a \geq 0$，成立：
\begin{equation*}
    \begin{aligned}
        T(v+a)(k) & =\max _x[u(f(k)-x)+\beta(v(x)+a)] \\
            & =\max _x[u(f(k)-x)+\beta v(x)]+\beta a \\
            & =T v(k)+\beta a
    \end{aligned}
\end{equation*}
因此，贝尔曼算子$T$满足Blackwell充分条件，从而是压缩映射。

再次回忆此前提到所面临的两个问题：
\begin{itemize}
    \item 理论层面，贝尔曼方程$v=Tv$有解吗，即是否存在唯一不动点
    \item 算法层面，构造的值函数序列是否收敛
\end{itemize}

在证明了贝尔曼算子$T$是压缩映射后，直接应用压缩映射定理，似乎我们已经完成了证明。但实际上，还存在两个确实缺陷：
一个存在于证明贝尔曼算子压缩性质（即运用Blackwell充分条件）的过程中；另一个则存在于运用压缩映射定理的过程中。具体看：

（1）\textbf{Blackwell充分条件应用的前提}：算子$T$是定义在\textbf{有界函数空间}上的。

对于递归形式下的最优增长模型，我们所分析的也是一个函数空间，
即$S$中的元素是值函数$v$。但是，值函数是否有界？答案是肯定的，这与我们对生产函数和效用函数的假设密切相关。
根据Inada条件、$\beta < 1$和消费与资本的非负性，由级数收敛性知识，下界的存在是显然的。
对于上界：由生产函数的凹性，资本积累存在上界，记为$k \leq k^{\star}$，从而$f\left(k\right) \leq f\left(k^{\star}\right)$；
消费的最大可能取值则为$f\left(k\right)$，即下期资本取0，从而有：
$u\left(c\right)\leq u\left(f\left(k\right)\right)\leq u\left(f\left(k^{\star}\right)\right)$。因此，
\begin{equation*}
    v\left(k\right) = \sum_{t=0}^{\infty} \beta^t u\left(c\right) \leq 
        \sum_{t=0}^{\infty} \beta^t u\left(f\left(k^{\star}\right)\right) = 
            \frac{u\left(f\left(k^{\star}\right)\right)}{1-\beta} < \infty
\end{equation*}
从而值函数有界，可以对贝尔曼算子应用Blackwell充分条件。

（2）\textbf{压缩映射定理应用的前提}：压缩映射$T$是定义在\textbf{完备空间}$\left(S,d\right)$上的。

值函数是有界的，但是值函数空间否是完备的呢？我们曾在\cref{th-complete-continuous-func-space}中指出，
定义了上确界范数$\|f\|=\sup _{x \in X}(|f(x)|)$的\textbf{有界连续函数集合}空间$C\left(X\right)$是一个完备空间。
那么，只要能证明\textbf{值函数是连续的}，我们就能在Blackwell条件的基础上进一步应用压缩映射定理：
证明不动点的存在性和唯一性，并且运用其推广证明值函数序列的收敛性。